\chapter{Types de médias multiples}
\label{ch:multimedia}

mybibli n'est pas réservé aux livres. Quatre types de
médias\index{types de médias} sont des citoyens de première classe :

\begin{itemize}
  \item \textbf{Livre} — le défaut. Romans, essais, ouvrages de
        référence, albums illustrés pour enfants.
  \item \textbf{BD / comic}\index{BD / comic} — comic individuel ou
        album omnibus multi-positions (par exemple une « anthologie
        Tintin » qui réunit physiquement les volumes 1 à 5).
  \item \textbf{Audio} — CD, vinyles, livres audio sur support
        physique.
  \item \textbf{Film / série}\index{film} — DVD, coffrets Blu-ray,
        intégrales de séries TV.
\end{itemize}

Le type de média pilote trois choses :

\begin{enumerate}
  \item \textbf{Quel fournisseur de métadonnées est consulté en
        premier.} Les ISBN routent vers BnF + Google Books + Open
        Library ; les films routent vers TMDb + OMDb ; l'audio
        route vers MusicBrainz.
  \item \textbf{Quels champs sont exposés sur la page du titre.}
        Les titres audio font apparaître une track list ; les
        films font apparaître réalisateur, durée, restriction
        d'âge.
  \item \textbf{L'icône affichée à côté du titre} dans les
        résultats de recherche et la liste de catalogue.
\end{enumerate}

\section{Livres}

Le type de média par défaut. Catalogage tel que décrit en
section~\ref{sec:catalog-scan}. Les séries multi-volumes
(encyclopédies, trilogies, \dots) sont modélisées par une entité
\emph{série}\index{série} au chapitre~\ref{ch:configuration}.

\section{BD et comics}

Un volume \textbf{omnibus}\index{omnibus} peut contenir plusieurs
positions dans une série. mybibli gère cela via la fonctionnalité
\emph{multi-position}\index{multi-position} : au moment de la
création d'un volume, vous pouvez fixer la \emph{plage de
positions} (par exemple « contient les volumes 5 à 7 »). La
détection de lacunes de série (chapitre~\ref{ch:usage}) comprend ces
plages et ne signalera pas les positions 5, 6 ou 7 comme manquantes
une fois l'omnibus en rayon.

\section{Sorties audio}

Les titres audio utilisent MusicBrainz\index{MusicBrainz} comme
source principale de métadonnées. Le code-barres au dos d'un CD
(EAN-13) se résout de la même façon qu'un ISBN.

\begin{itemize}
  \item \textbf{Pistes.} MusicBrainz renvoie la liste des pistes ;
        mybibli les affiche sur la page du titre.
  \item \textbf{Support.} CD vs vinyle vs cassette est capturé sur
        la ligne du volume (pas sur le titre) — deux pressages
        identiques sur supports différents sont deux volumes du
        même titre.
\end{itemize}

\section{Films et séries}

Les films et séries TV utilisent OMDb\index{OMDb} et
TMDb\index{TMDb} comme sources principales de métadonnées. Un
coffret de série est modélisé comme un \emph{titre} avec un volume
multi-position couvrant la plage de disques.

\begin{tip}
OMDb est gratuit pour de faibles volumes (moins de 1000 requêtes
par jour) ; TMDb est gratuit avec attribution. L'un ou l'autre
couvre largement une étagère DVD / Blu-ray familiale. Voir le
chapitre~\ref{ch:providers} pour la mise en place.
\end{tip}

\section{Choisir le type de média au scan}

Après scan d'un code-barres inconnu, mybibli devine le type de
média depuis le préfixe EAN (le préfixe pays GS1 et la réponse du
fournisseur). La devinette est affichée sur la page de catalogage ;
vous pouvez la corriger avant de cliquer \textbf{Ajouter un volume}.
Re-typer un titre plus tard se fait en quelques clics
(\texttt{/title/<id>}) et la chaîne de fournisseurs ré-exécute
contre le nouveau type si vous déclenchez \emph{Re-fetch
metadata}.
